\section{Counterexamples on thirteen and fourteen vertices}\label{sec:examples}

The two assertions in Conjecture~14 are refuted at different orders.
Odd integral torsion supplies a graph with trivial homology over
$\Ztwo$ outside $\IS$. A second construction belongs to $\I$ but admits
no first strong deletion. Adding one edge to that construction also
gives a counterexample to Conjecture~15.

Throughout this section, $\Gt$ has vertex set $\{0,\ldots,12\}$ and
$\Ga,\Gb,\Gc$ have vertex set $\{0,\ldots,13\}$. Their edges are defined
by Table~\ref{tab:witness-adjacency}. The figures use these same labels.

\subsection{Odd torsion on thirteen vertices}\label{sec:thirteen}

The clique complex of $\Gt$ is obtained from a triangulated disc by
wrapping its boundary three times around a four-cycle. This description accounts
for the torsion and gives a direct proof of the first refutation.

\begin{theorem}\label{thm:A}
Let $\Gt'=\Gt+\{4,7\}$ and $\Gt''=\Gt'+\{5,6\}$.
The graphs $\Gt,\Gt',\Gt''$ have $40$, $41$ and $42$ edges, respectively,
and their clique complexes are homotopy equivalent to $M(\ZZ/3,1)$.
In particular, they are $\Ztwo$-acyclic and none belongs to $\I$.
They refute the first assertion of Conjecture~14.
\end{theorem}

\begin{proof}
The graph $\Gt$ is $K_4$-free and has $28$ triangles. Each of its edges
lies in two triangles, except for the four edges of
$8\text{--}10\text{--}9\text{--}11\text{--}8$, which lie in three.
Cutting along that cycle gives the triangulated disc on the left of
Figure~\ref{fig:G13}. An explicit triangle list and an order of attachment
proving that it is a disc are given in
Appendix~\ref{app:disc-certificates}.
Its boundary labels, starting at the top and read clockwise, are
\[
 10,9,11,8\mid10,9,11,8\mid10,9,11,8.
\]
Each interior vertex occurs once. Identifying equally labelled boundary
vertices and their corresponding boundary edges recovers precisely
$\cl(\Gt)$.
Proposition~\ref{prop:disc}, with attaching degree $3$, therefore gives
\[
 \cl(\Gt)\simeq M(\ZZ/3,1),\qquad
 \widetilde H_1(\Gt;\ZZ)\cong\ZZ/3,
 \qquad \widetilde H_i(\Gt;\ZZ)=0\quad(i\ne1).
\]
The universal coefficient theorem gives vanishing reduced homology
over $\Ztwo$ in every degree.

For each of the two specified edge additions, the common neighbourhood
at the time of addition is the singleton $\{12\}$. Both additions are
admissible and preserve the homotopy type of the clique complex.
Nonzero integral homology excludes all three graphs from $\I$, hence
also from $\IS$. Since they have fewer than the $15$ vertices of the
graph in Conjecture~14, they refute its first assertion.
\end{proof}

\begin{figure}[tbp]
\centering
\begin{tikzpicture}[x=1cm,y=1cm]
\definecolor{witnessred}{RGB}{181,55,36}
\definecolor{witnessblue}{RGB}{32,89,140}
\coordinate (v0) at (0.942769700,-1.632925019);
\coordinate (v1) at (1.632925019,0.942769700);
\coordinate (v2) at (-1.632925019,-0.942769700);
\coordinate (v3) at (-0.942769700,1.632925019);
\coordinate (v4) at (1.981676324,-0.530988571);
\coordinate (v5) at (-0.530988571,-1.981676324);
\coordinate (v6) at (0.530988571,1.981676324);
\coordinate (v7) at (-1.981676324,0.530988571);
\coordinate (v8) at (-1.700000000,2.944486373);
\coordinate (v9) at (3.400000000,0.000000000);
\coordinate (v10) at (-1.700000000,-2.944486373);
\coordinate (v11) at (1.700000000,-2.944486373);
\coordinate (v12) at (1.700000000,2.944486373);
\coordinate (v13) at (-3.400000000,-0.000000000);
\coordinate (v14) at (-2.944486373,-1.700000000);
\coordinate (v15) at (2.944486373,-1.700000000);
\coordinate (v16) at (0.000000000,3.400000000);
\coordinate (v17) at (2.944486373,1.700000000);
\coordinate (v18) at (0.000000000,-3.400000000);
\coordinate (v19) at (-2.944486373,1.700000000);
\coordinate (v20) at (0.000000000,0.000000000);
\fill[blue!3] (v0) -- (v4) -- (v15) -- cycle;
\fill[blue!3] (v0) -- (v4) -- (v20) -- cycle;
\fill[blue!3] (v0) -- (v5) -- (v18) -- cycle;
\fill[blue!3] (v0) -- (v5) -- (v20) -- cycle;
\fill[blue!3] (v0) -- (v11) -- (v15) -- cycle;
\fill[blue!3] (v0) -- (v11) -- (v18) -- cycle;
\fill[blue!3] (v1) -- (v4) -- (v9) -- cycle;
\fill[blue!3] (v1) -- (v4) -- (v20) -- cycle;
\fill[blue!3] (v1) -- (v6) -- (v12) -- cycle;
\fill[blue!3] (v1) -- (v6) -- (v20) -- cycle;
\fill[blue!3] (v1) -- (v9) -- (v17) -- cycle;
\fill[blue!3] (v1) -- (v12) -- (v17) -- cycle;
\fill[blue!3] (v2) -- (v5) -- (v10) -- cycle;
\fill[blue!3] (v2) -- (v5) -- (v20) -- cycle;
\fill[blue!3] (v2) -- (v7) -- (v13) -- cycle;
\fill[blue!3] (v2) -- (v7) -- (v20) -- cycle;
\fill[blue!3] (v2) -- (v10) -- (v14) -- cycle;
\fill[blue!3] (v2) -- (v13) -- (v14) -- cycle;
\fill[blue!3] (v3) -- (v6) -- (v16) -- cycle;
\fill[blue!3] (v3) -- (v6) -- (v20) -- cycle;
\fill[blue!3] (v3) -- (v7) -- (v19) -- cycle;
\fill[blue!3] (v3) -- (v7) -- (v20) -- cycle;
\fill[blue!3] (v3) -- (v8) -- (v16) -- cycle;
\fill[blue!3] (v3) -- (v8) -- (v19) -- cycle;
\fill[blue!3] (v4) -- (v9) -- (v15) -- cycle;
\fill[blue!3] (v5) -- (v10) -- (v18) -- cycle;
\fill[blue!3] (v6) -- (v12) -- (v16) -- cycle;
\fill[blue!3] (v7) -- (v13) -- (v19) -- cycle;
\draw[draw=black!52,line width=.38pt] (v0) -- (v4);
\draw[draw=black!52,line width=.38pt] (v0) -- (v5);
\draw[draw=black!52,line width=.38pt] (v0) -- (v11);
\draw[draw=black!52,line width=.38pt] (v0) -- (v15);
\draw[draw=black!52,line width=.38pt] (v0) -- (v18);
\draw[draw=black!52,line width=.38pt] (v0) -- (v20);
\draw[draw=black!52,line width=.38pt] (v1) -- (v4);
\draw[draw=black!52,line width=.38pt] (v1) -- (v6);
\draw[draw=black!52,line width=.38pt] (v1) -- (v9);
\draw[draw=black!52,line width=.38pt] (v1) -- (v12);
\draw[draw=black!52,line width=.38pt] (v1) -- (v17);
\draw[draw=black!52,line width=.38pt] (v1) -- (v20);
\draw[draw=black!52,line width=.38pt] (v2) -- (v5);
\draw[draw=black!52,line width=.38pt] (v2) -- (v7);
\draw[draw=black!52,line width=.38pt] (v2) -- (v10);
\draw[draw=black!52,line width=.38pt] (v2) -- (v13);
\draw[draw=black!52,line width=.38pt] (v2) -- (v14);
\draw[draw=black!52,line width=.38pt] (v2) -- (v20);
\draw[draw=black!52,line width=.38pt] (v3) -- (v6);
\draw[draw=black!52,line width=.38pt] (v3) -- (v7);
\draw[draw=black!52,line width=.38pt] (v3) -- (v8);
\draw[draw=black!52,line width=.38pt] (v3) -- (v16);
\draw[draw=black!52,line width=.38pt] (v3) -- (v19);
\draw[draw=black!52,line width=.38pt] (v3) -- (v20);
\draw[draw=black!52,line width=.38pt] (v4) -- (v9);
\draw[draw=black!52,line width=.38pt] (v4) -- (v15);
\draw[draw=black!52,line width=.38pt] (v4) -- (v20);
\draw[draw=black!52,line width=.38pt] (v5) -- (v10);
\draw[draw=black!52,line width=.38pt] (v5) -- (v18);
\draw[draw=black!52,line width=.38pt] (v5) -- (v20);
\draw[draw=black!52,line width=.38pt] (v6) -- (v12);
\draw[draw=black!52,line width=.38pt] (v6) -- (v16);
\draw[draw=black!52,line width=.38pt] (v6) -- (v20);
\draw[draw=black!52,line width=.38pt] (v7) -- (v13);
\draw[draw=black!52,line width=.38pt] (v7) -- (v19);
\draw[draw=black!52,line width=.38pt] (v7) -- (v20);
\draw[draw=witnessred,line width=.8pt] (v8) -- (v16);
\draw[draw=witnessred,line width=.8pt] (v8) -- (v19);
\draw[draw=witnessred,line width=.8pt] (v9) -- (v15);
\draw[draw=witnessred,line width=.8pt] (v9) -- (v17);
\draw[draw=witnessred,line width=.8pt] (v10) -- (v14);
\draw[draw=witnessred,line width=.8pt] (v10) -- (v18);
\draw[draw=witnessred,line width=.8pt] (v11) -- (v15);
\draw[draw=witnessred,line width=.8pt] (v11) -- (v18);
\draw[draw=witnessred,line width=.8pt] (v12) -- (v16);
\draw[draw=witnessred,line width=.8pt] (v12) -- (v17);
\draw[draw=witnessred,line width=.8pt] (v13) -- (v14);
\draw[draw=witnessred,line width=.8pt] (v13) -- (v19);
\node[jfef vertex] at (v0) {$0$};
\node[jfef vertex] at (v1) {$1$};
\node[jfef vertex] at (v2) {$2$};
\node[jfef vertex] at (v3) {$3$};
\node[jfef vertex] at (v4) {$4$};
\node[jfef vertex] at (v5) {$5$};
\node[jfef vertex] at (v6) {$6$};
\node[jfef vertex] at (v7) {$7$};
\node[jfef vertex,draw=witnessred,line width=.7pt] at (v8) {$8$};
\node[jfef vertex,draw=witnessred,line width=.7pt] at (v9) {$8$};
\node[jfef vertex,draw=witnessred,line width=.7pt] at (v10) {$8$};
\node[jfef vertex,draw=witnessred,line width=.7pt] at (v11) {$9$};
\node[jfef vertex,draw=witnessred,line width=.7pt] at (v12) {$9$};
\node[jfef vertex,draw=witnessred,line width=.7pt] at (v13) {$9$};
\node[jfef vertex,draw=witnessred,line width=.7pt] at (v14) {$10$};
\node[jfef vertex,draw=witnessred,line width=.7pt] at (v15) {$10$};
\node[jfef vertex,draw=witnessred,line width=.7pt] at (v16) {$10$};
\node[jfef vertex,draw=witnessred,line width=.7pt] at (v17) {$11$};
\node[jfef vertex,draw=witnessred,line width=.7pt] at (v18) {$11$};
\node[jfef vertex,draw=witnessred,line width=.7pt] at (v19) {$11$};
\node[jfef vertex] at (v20) {$12$};
\end{tikzpicture}
\hspace{6mm}\begin{tikzpicture}[x=1cm,y=1cm]
\definecolor{witnessred}{RGB}{181,55,36}
\definecolor{witnessblue}{RGB}{32,89,140}
\coordinate (v0) at (1.146661411,2.139452559);
\coordinate (v1) at (-2.265488516,1.039851549);
\coordinate (v2) at (2.265489343,-1.039852163);
\coordinate (v3) at (-1.146663831,-2.139454358);
\coordinate (v4) at (-1.146661411,2.139452559);
\coordinate (v5) at (2.265488516,1.039851549);
\coordinate (v6) at (-2.265489343,-1.039852163);
\coordinate (v7) at (1.146663831,-2.139454358);
\coordinate (v8) at (3.440000000,-3.200000000);
\coordinate (v9) at (-3.440000000,3.200000000);
\coordinate (v10) at (3.440000000,3.200000000);
\coordinate (v11) at (-3.440000000,-3.200000000);
\coordinate (v12) at (0.000000000,0.000000000);
\draw[draw=black!52,line width=.38pt] (v0) -- (v4);
\draw[draw=black!52,line width=.38pt] (v0) -- (v5);
\draw[draw=black!52,line width=.38pt] (v0) -- (v9);
\draw[draw=black!52,line width=.38pt] (v0) -- (v10);
\draw[draw=black!52,line width=.38pt] (v0) -- (v11);
\draw[draw=black!52,line width=.38pt] (v0) -- (v12);
\draw[draw=black!52,line width=.38pt] (v1) -- (v4);
\draw[draw=black!52,line width=.38pt] (v1) -- (v6);
\draw[draw=black!52,line width=.38pt] (v1) -- (v8);
\draw[draw=black!52,line width=.38pt] (v1) -- (v9);
\draw[draw=black!52,line width=.38pt] (v1) -- (v11);
\draw[draw=black!52,line width=.38pt] (v1) -- (v12);
\draw[draw=black!52,line width=.38pt] (v2) -- (v5);
\draw[draw=black!52,line width=.38pt] (v2) -- (v7);
\draw[draw=black!52,line width=.38pt] (v2) -- (v8);
\draw[draw=black!52,line width=.38pt] (v2) -- (v9);
\draw[draw=black!52,line width=.38pt] (v2) -- (v10);
\draw[draw=black!52,line width=.38pt] (v2) -- (v12);
\draw[draw=black!52,line width=.38pt] (v3) -- (v6);
\draw[draw=black!52,line width=.38pt] (v3) -- (v7);
\draw[draw=black!52,line width=.38pt] (v3) -- (v8);
\draw[draw=black!52,line width=.38pt] (v3) -- (v10);
\draw[draw=black!52,line width=.38pt] (v3) -- (v11);
\draw[draw=black!52,line width=.38pt] (v3) -- (v12);
\draw[draw=black!52,line width=.38pt] (v4) -- (v8);
\draw[draw=black!52,line width=.38pt] (v4) -- (v10);
\draw[draw=black!52,line width=.38pt] (v4) -- (v12);
\draw[draw=black!52,line width=.38pt] (v5) -- (v8);
\draw[draw=black!52,line width=.38pt] (v5) -- (v11);
\draw[draw=black!52,line width=.38pt] (v5) -- (v12);
\draw[draw=black!52,line width=.38pt] (v6) -- (v9);
\draw[draw=black!52,line width=.38pt] (v6) -- (v10);
\draw[draw=black!52,line width=.38pt] (v6) -- (v12);
\draw[draw=black!52,line width=.38pt] (v7) -- (v9);
\draw[draw=black!52,line width=.38pt] (v7) -- (v11);
\draw[draw=black!52,line width=.38pt] (v7) -- (v12);
\draw[draw=witnessred,line width=.8pt] (v8) -- (v10);
\draw[draw=witnessred,line width=.8pt] (v8) -- (v11);
\draw[draw=witnessred,line width=.8pt] (v9) -- (v10);
\draw[draw=witnessred,line width=.8pt] (v9) -- (v11);
\node[jfef vertex] at (v0) {$0$};
\node[jfef vertex] at (v1) {$1$};
\node[jfef vertex] at (v2) {$2$};
\node[jfef vertex] at (v3) {$3$};
\node[jfef vertex] at (v4) {$4$};
\node[jfef vertex] at (v5) {$5$};
\node[jfef vertex] at (v6) {$6$};
\node[jfef vertex] at (v7) {$7$};
\node[jfef vertex,draw=witnessred,line width=.7pt] at (v8) {$8$};
\node[jfef vertex,draw=witnessred,line width=.7pt] at (v9) {$9$};
\node[jfef vertex,draw=witnessred,line width=.7pt] at (v10) {$10$};
\node[jfef vertex,draw=witnessred,line width=.7pt] at (v11) {$11$};
\node[jfef vertex] at (v12) {$12$};
\end{tikzpicture}

\caption{On the left, a triangulated disc whose equally labelled boundary
vertices and corresponding boundary edges are identified to obtain
$\cl(\Gt)$. The boundary traverses
$10$--$9$--$11$--$8$ three times. The $28$ small triangular regions are
its faces. On the right, the graph $\Gt$; the four edges in three
triangles are highlighted. The highlighted boundary on the left maps
to the highlighted four-cycle on the right.}\label{fig:G13}
\end{figure}

The construction also attains the least possible order for odd torsion
in the first homology of a flag complex. This conclusion uses the
earlier classification of torsion graphs, without an acyclicity
assumption.

\begin{proposition}\label{rem:adamaszek}
The minimum number of vertices of a flag complex whose first integral
homology has odd torsion is $13$.
\end{proposition}

\begin{proof}
A torsion graph is called irreducible if every vertex deletion removes
all torsion from its first homology. Adamaszek's classification
\cite[Theorem~1.1 and Section~4]{Adamaszek2014} gives no torsion through
ten vertices, four torsion graphs on eleven vertices, each with first
homology $\ZZ/2$, and $363$ irreducible torsion graphs on twelve vertices,
each with torsion subgroup $\ZZ/2$. The absence of torsion through ten vertices and the four torsion
graphs on eleven vertices also follow from Katzman
\cite[Theorem~4.1 and Section~5]{Katzman2006}.

If a twelve-vertex torsion graph $G$ is not irreducible, it has a vertex
$v$ such that $G-v$ is one of the four eleven-vertex graphs. If $v$ is
isolated, then $H_1(\cl(G);\ZZ)=H_1(\cl(G-v);\ZZ)\cong\ZZ/2$.
Otherwise, apply the Mayer--Vietoris sequence to
\[
 \cl(G)=\cl(G-v)\cup\bigl(v*\cl(\Nb_G(v))\bigr).
\]
It gives an exact sequence
$0\to Q\to H_1(\cl(G);\ZZ)\to F\to0$, where $Q$ is a quotient of
$\ZZ/2$ and $F$ is a subgroup of
$\widetilde H_0(\cl(\Nb_G(v));\ZZ)$. The group $F$ is free abelian,
so the sequence splits. Thus $H_1(\cl(G);\ZZ)$ has no odd torsion.
Theorem~\ref{thm:A} attains the lower bound.
\end{proof}

\subsection{A non-strong I-contractible graph on fourteen vertices}
\label{sec:fourteen}

For the second assertion of Conjecture~14, nonzero homology cannot
provide the obstruction: an I-contractible graph is integrally
acyclic. Instead, we show that every possible first strong deletion
is obstructed, and then establish I-contractibility by an edge addition.

\begin{theorem}\label{thm:B}
The graph $\Ga$ is $K_4$-free, has $14$ vertices and $44$ edges, and
belongs to $\I\setminus\IS$. Its one-edge extension $\Ga+\{0,1\}$ belongs to
$\IS$, and deletion of that edge is $\IS$-admissible.
Consequently, $\Ga$ refutes the second assertion of Conjecture~14,
and $\Ga+\{0,1\}$ refutes Conjecture~15.
\end{theorem}

\begin{proof}
The common neighbourhood of $0$ and $1$ in $\Ga$ is the induced path
\[
 8\text{--}12\text{--}5\text{--}10\text{--}13.
\]
Since nonempty trees belong to $\IS$ (Section~\ref{sec:contractible}),
adding $\{0,1\}$ is admissible.
Starting from $\Ga+\{0,1\}$, perform the following deletions in the
displayed order:
\[
\begin{gathered}
 e(0,5),\ e(0,10),\ e(0,12),\ e(8,12),\ e(1,12),\\
 v(5),\ v(10),\ v(3),\ v(6),\ v(2),\ v(7),\ v(1),\ v(0),\\
 v(4),\ v(8),\ v(11),\ v(12),\ v(9).
\end{gathered}
\]
Here $e(u,w)$ deletes an edge and $v(u)$ deletes a vertex; labels do
not change after a deletion. Every tested neighbourhood is a nonempty
tree, as the explicit lists in Table~\ref{tab:Ga-reduction} show.
The final graph is the singleton $13$. Hence $\Ga+\{0,1\}\in\IS$.
Deleting the added edge is an I-admissible transformation, with the
same common neighbourhood as before its addition, so $\Ga\in\I$.

To prove $\Ga\notin\IS$, consider its links. Each vertex neighbourhood
is a connected triangle-free graph; its first Betti number is positive
in every case, as recorded in Table~\ref{tab:witness-links}.
Every edge has two or three common neighbours, and those neighbours
form an independent set because $\Ga$ is $K_4$-free.
Thus neither a vertex neighbourhood nor an edge neighbourhood is
acyclic. None belongs to $\IS$, so no first strong deletion is possible.

The two graphs have $14$ vertices, whereas both examples conjectured
to be smallest in the earlier paper have $15$ vertices. This proves
the stated refutations. The global lower bounds are established from
the complete census below.
\end{proof}

\begin{figure}[tbp]
\centering
\begin{tikzpicture}[x=1cm,y=1cm]
\definecolor{witnessred}{RGB}{181,55,36}
\definecolor{witnessblue}{RGB}{32,89,140}
\coordinate (v0) at (4.899000000,-0.124598525);
\coordinate (v1) at (2.722969139,3.944383940);
\coordinate (v2) at (-2.435697872,0.406103749);
\coordinate (v3) at (-4.899000000,4.045228676);
\coordinate (v4) at (1.271029059,-0.654558239);
\coordinate (v5) at (4.387943889,-3.466555902);
\coordinate (v6) at (-4.899000000,-4.554000000);
\coordinate (v7) at (-2.939555778,-3.846959139);
\coordinate (v8) at (-0.868699866,3.253825943);
\coordinate (v9) at (2.353335221,-2.606664647);
\coordinate (v10) at (3.219437385,-4.554000000);
\coordinate (v11) at (3.996206901,3.171731296);
\coordinate (v12) at (1.198003432,4.416000000);
\coordinate (v13) at (-1.343328314,-3.846959139);
\draw[draw=witnessblue,line width=1pt,dashed] (v0) -- (v1);
\draw[draw=black!52,line width=.38pt] (v0) -- (v4);
\draw[draw=black!52,line width=.38pt] (v0) -- (v5);
\draw[draw=black!52,line width=.38pt] (v0) -- (v8);
\draw[draw=black!52,line width=.38pt] (v0) -- (v10);
\draw[draw=black!52,line width=.38pt] (v0) -- (v12);
\draw[draw=black!52,line width=.38pt] (v0) -- (v13);
\draw[draw=black!52,line width=.38pt] (v1) -- (v5);
\draw[draw=black!52,line width=.38pt] (v1) -- (v7);
\draw[draw=black!52,line width=.38pt] (v1) -- (v8);
\draw[draw=black!52,line width=.38pt] (v1) -- (v10);
\draw[draw=black!52,line width=.38pt] (v1) -- (v12);
\draw[draw=black!52,line width=.38pt] (v1) -- (v13);
\draw[draw=black!52,line width=.38pt] (v2) -- (v6);
\draw[draw=black!52,line width=.38pt] (v2) -- (v7);
\draw[draw=black!52,line width=.38pt] (v2) -- (v9);
\draw[draw=black!52,line width=.38pt] (v2) -- (v12);
\draw[draw=black!52,line width=.38pt] (v2) -- (v13);
\draw[draw=black!52,line width=.38pt] (v3) -- (v6);
\draw[draw=black!52,line width=.38pt] (v3) -- (v9);
\draw[draw=black!52,line width=.38pt] (v3) -- (v10);
\draw[draw=black!52,line width=.38pt] (v3) -- (v12);
\draw[draw=black!52,line width=.38pt] (v3) -- (v13);
\draw[draw=black!52,line width=.38pt] (v4) -- (v8);
\draw[draw=black!52,line width=.38pt] (v4) -- (v9);
\draw[draw=black!52,line width=.38pt] (v4) -- (v11);
\draw[draw=black!52,line width=.38pt] (v4) -- (v13);
\draw[draw=black!52,line width=.38pt] (v5) -- (v9);
\draw[draw=witnessred,line width=.8pt] (v5) -- (v10);
\draw[draw=black!52,line width=.38pt] (v5) -- (v11);
\draw[draw=witnessred,line width=.8pt] (v5) -- (v12);
\draw[draw=black!52,line width=.38pt] (v6) -- (v7);
\draw[draw=black!52,line width=.38pt] (v6) -- (v10);
\draw[draw=black!52,line width=.38pt] (v6) -- (v11);
\draw[draw=black!52,line width=.38pt] (v6) -- (v12);
\draw[draw=black!52,line width=.38pt] (v7) -- (v8);
\draw[draw=black!52,line width=.38pt] (v7) -- (v11);
\draw[draw=black!52,line width=.38pt] (v7) -- (v13);
\draw[draw=black!52,line width=.38pt] (v8) -- (v11);
\draw[draw=witnessred,line width=.8pt] (v8) -- (v12);
\draw[draw=black!52,line width=.38pt] (v9) -- (v11);
\draw[draw=black!52,line width=.38pt] (v9) -- (v12);
\draw[draw=black!52,line width=.38pt] (v9) -- (v13);
\draw[draw=black!52,line width=.38pt] (v10) -- (v11);
\draw[draw=witnessred,line width=.8pt] (v10) -- (v13);
\node[jfef vertex,draw=witnessblue,line width=.8pt] at (v0) {$0$};
\node[jfef vertex,draw=witnessblue,line width=.8pt] at (v1) {$1$};
\node[jfef vertex] at (v2) {$2$};
\node[jfef vertex] at (v3) {$3$};
\node[jfef vertex] at (v4) {$4$};
\node[jfef vertex,fill=orange!12] at (v5) {$5$};
\node[jfef vertex] at (v6) {$6$};
\node[jfef vertex] at (v7) {$7$};
\node[jfef vertex,fill=orange!12] at (v8) {$8$};
\node[jfef vertex] at (v9) {$9$};
\node[jfef vertex,fill=orange!12] at (v10) {$10$};
\node[jfef vertex] at (v11) {$11$};
\node[jfef vertex,fill=orange!12] at (v12) {$12$};
\node[jfef vertex,fill=orange!12] at (v13) {$13$};
\end{tikzpicture}

\caption{The graph $\Ga$ consists of the solid edges. Adding the dashed
edge $\{0,1\}$ gives a strong I-contractible graph. The common
neighbourhood of its endpoints is the highlighted induced path
$8$--$12$--$5$--$10$--$13$. Deleting the dashed edge is admissible but
leaves a graph outside $\IS$.}\label{fig:G14}
\end{figure}

The order-$14$ classification also uses $\Gb$ and $\Gc$, two further
members of $\I\setminus\IS$ that become strong I-contractible after
one admissible edge addition. Table~\ref{tab:witness-adjacency} defines
them and Table~\ref{tab:variant-reductions} supplies their reductions.

\begin{proposition}\label{prop:other-seeds}
The graphs $\Gb$ and $\Gc$ are $K_4$-free members of
$\I\setminus\IS$, each with $14$ vertices and $45$ edges.
The extensions $\Gb+\{0,3\}$ and $\Gc+\{2,3\}$ belong to $\IS$,
and the added edges are admissible.
\end{proposition}

\begin{proof}
The common neighbourhoods for the two additions are the induced paths
$7\text{--}13\text{--}9\text{--}12$ and
$6\text{--}13\text{--}7\text{--}10\text{--}11$, respectively.
In each column of Table~\ref{tab:variant-reductions}, every tested
neighbourhood is a nonempty tree and the last vertex is $13$.
This proves membership of both extensions in $\IS$ and of the original
graphs in $\I$. Table~\ref{tab:witness-links} excludes every possible
first strong deletion from $\Gb$ and $\Gc$, just as for $\Ga$.
The graph $\Gb$ has an edge with four common neighbours, whereas
$\Gc$ has none. Together with their edge counts, this shows that
$\Ga,\Gb,\Gc$ are pairwise non-isomorphic.
\end{proof}

The same local structure also obstructs simplicial collapses.

\begin{proposition}\label{prop:noncollapsible}
The clique complexes of $\Ga,\Gb,\Gc$ are contractible and
non-collapsible.
\end{proposition}

\begin{proof}
Their I-contractibility implies contractibility of their clique
complexes. Each complex is two-dimensional. By
Table~\ref{tab:witness-links}, every edge lies in at least two triangles,
and every vertex lies in a triangle. There is therefore no free face
with which to start an elementary collapse.
\end{proof}
